% Generated draft: requires mathematical and source review.
\documentclass[11pt]{article}
\usepackage[margin=25mm]{geometry}
\usepackage{fontspec}
\setmainfont{Latin Modern Roman}
\usepackage{amsmath,amssymb,mathtools,unicode-math}
\setmathfont{Latin Modern Math}
\usepackage{enumitem,needspace,xurl,hyperref}
\setlength{\parindent}{0pt}
\setlength{\parskip}{5pt}
\setlength{\emergencystretch}{3em}
\newcommand{\N}{\mathbb{N}}
\newcommand{\Z}{\mathbb{Z}}
\newcommand{\Q}{\mathbb{Q}}
\newcommand{\R}{\mathbb{R}}
\newcommand{\C}{\mathbb{C}}
\newcommand{\F}{\mathbb{F}}
\newcommand{\A}{\mathbb{A}}
\newcommand{\PP}{\mathbb{P}}
\newcommand{\E}{\mathbb{E}}
\newcommand{\eps}{\varepsilon}
\newcommand{\dd}{\,\mathrm{d}}
\begin{document}
\section*{UnsolvedMath ID: 800070}
\textbf{Source problem code:} catalogue entry
\subsection{Definitions and mathematical statement}
In the standard setting, let \(M\) be a smooth four-dimensional manifold. An Engel structure on \(M\) is a smooth rank-two distribution \(\mathcal D\subset TM\) such that
\[
\mathcal E:=\mathcal D+[\mathcal D,\mathcal D],\qquad
\operatorname{rank}(\mathcal E)=3,\qquad
\mathcal E+[\mathcal D,\mathcal E]=TM.
\]
Here brackets of distributions mean the pointwise spans of Lie brackets of their local smooth sections. The characteristic, or kernel, line field \(\mathcal W\subset\mathcal D\) is the kernel of the Levi form
\[
L_{\mathcal E,p}\colon\Lambda^2\mathcal E_p\longrightarrow T_pM/\mathcal E_p,
\qquad
L_{\mathcal E,p}(u,v)=[U,V]_p\bmod\mathcal E_p,
\]
where \(U,V\) are local sections of \(\mathcal E\) extending \(u,v\). Thus
\[
\mathcal W_p=\{u\in\mathcal E_p:L_{\mathcal E,p}(u,v)=0\text{ for every }v\in\mathcal E_p\}.
\]
Its integral curves define a one-dimensional foliation. The source does not specify an orientation or a parametrization of these curves, so a particular generating flow is not part of the supplied data.

\textbf{Source question.} Are there any dynamical properties of \(\mathcal W\) that persist under homotopies of the Engel structure?

For a family \((\mathcal D_t)_{t\in[0,1]}\) of Engel structures on a fixed \(M\), write \(\mathcal W_t\) for the corresponding kernel line fields. One possible meaning of persistence for a property \(P\) is the requirement
\[
P(\mathcal W_0)\ \Longrightarrow\ P(\mathcal W_t)
\quad\text{for every }t\in[0,1]
\]
for every allowed homotopy. This is only a possible formal interpretation: the source explicitly calls its question vague and does not specify the admissible properties, the class of manifolds, the regularity or boundary conditions of homotopies, or whether persistence means invariance under arbitrary homotopies or robustness under small perturbations. Consequently, the supplied record does not determine a unique precise conjecture.

The source attributes some constraints on the dynamics of \(\mathcal W\) to Nicola Pia, while saying that little is known overall. It also states that Engel structures lack stability in the sense that perturbations can yield Engel structures that need not be diffeomorphic to the original structure. These are supplied contextual qualifications, not independently verified assertions here.
\subsection{Short English statement}
For an Engel plane distribution on a smooth four-dimensional manifold, its associated rank-three distribution has a distinguished kernel line field. The question asks which dynamical features of the resulting one-dimensional foliation survive deformations through Engel structures. The source leaves the meaning of persistence and the allowed manifolds and deformations unspecified, so this remains a qualified research question rather than a uniquely determined conjecture.
\subsection{Sources}
\begin{itemize}
\item Local UnsolvedMath registry: supplied complete source record for permanent ID \(800070\), source problem code \texttt{catalogue entry}, titled \emph{Foliations and Diffeomorphism Groups --- Question 11}. The record attributes the problem to Nariman and Yazdi, with proposed year 2025. Its \texttt{open} label is supplied registry metadata, not an independently verified mathematical status.
\item Nariman and Yazdi, \emph{Problem List on Foliations and Diffeomorphism Groups} (2025), Question 11: \url{https://arxiv.org/abs/2503.18274}. The supplied provenance records access on 2026-10-16 and extraction by \texttt{arxiv\_latex}; no independent source visit or verification is claimed here.
\item The supplied question mentions work of Nicola Pia concerning constraints on kernel-line-field dynamics, but gives no title, theorem, bibliographic citation, or link for that work.
\item Registry provenance: upstream ID \(800070\), upstream revision \texttt{d371ee917ac9647fbb06cd6fc529deec6f3dda82}; record creation and update timestamps are both \texttt{2026-10-16T14:00:00Z}.
\end{itemize}
\end{document}
